% chapter1-body.tex — pure content, no preamble.
% This file is \input{} by both main.tex (book) and chapter1.tex (article).




\section{Basics}
In 2D,

\[
    \boxed{
        \begin{aligned}
            x^a & =(x,y),      & a & =1,2, \\
            q^i & =(r,\theta), & i & =1,2.
        \end{aligned}
    }
\]

\[
    \boxed{\partial_a=\frac{\partial}{\partial x^a}}
\]
\[
    \boxed{\partial_i=\frac{\partial}{\partial q^i}}
\]

The dual is

\[
    \boxed{
        \begin{aligned}
            dx^a
        \end{aligned}
    }
\]
\[
    \boxed{
        \begin{aligned}
            dq^i
        \end{aligned}
    }
\]

the transformation jacobian is

\[
    \partial_i x^a
\]
\[
    \partial_a q^i
\]

with jacobian we can transform as follow
\[
    \boxed{
        \begin{aligned}
            dx^a & =(\partial_i x^a)dq^i, \\
            dq^i & =(\partial_a q^i)dx^a.
        \end{aligned}
    }
\]

as well as partials

\[
    \boxed{
        \begin{aligned}
            \partial_i & =\partial_i x^a\partial_a, \\
            \partial_a & =\partial_a q^i\partial_i.
        \end{aligned}
    }
\]


\[
    v=v^a\partial_a
\]
\[
    \begin{aligned}
        \partial_a v^a
         & =(\partial_aq^i)\partial_i v^a                                                 \\
         & =(\partial_aq^i)\partial_i[(\partial_jx^a)v^j].                                \\
         & =(\partial_aq^i)[(\partial_i\partial_jx^a)v^j+(\partial_jx^a)(\partial_iv^j)]. \\
    \end{aligned}
\]



\[
    \partial_a v^a
    =
    \partial_i v^i
    +
    (\partial_aq^i)(\partial_i\partial_jx^a)v^j.
\]


\subsection{Propagation Directions in One Dimension}
In one dimension, the unit sphere consists of only two points,
\[
    S^0=\{-1,+1\},
\]
corresponding to the two possible propagation directions.

If
\[
    D^1=[-1,1],
\]
then
\[
    \partial D^1=S^0=\{-1,+1\}.
\]

For a direction $d\in S^0$, a plane-wave solution can be written as
\[
    u(x;d)=e^{ikxd}.
\]
Therefore,
\[
    d=+1 \quad\Longrightarrow\quad e^{ikx},
\]
and
\[
    d=-1 \quad\Longrightarrow\quad e^{-ikx}.
\]

% TODO: explain how the discrete directions in $S^0$
% generalize to $d\in S^{n-1}$ in higher dimensions.




\section{TMP}


\subsection{MVHR}
$$
    \boxed{
        \left[
            \begin{array}{c}
                \left[
                    \begin{array}{ccc}
                        [1\ 2\ 3\ 4]_{b} &
                        [5\ 6\ 7\ 8]_{b} &
                        [9\ 10\ 11\ 12]_{b}
                    \end{array}
                    \right]_{a}
                \\[12pt]
                \left[
                    \begin{array}{ccc}
                        [13\ 14\ 15\ 16]_{b} &
                        [17\ 18\ 19\ 20]_{b} &
                        [21\ 22\ 23\ 24]_{b}
                    \end{array}
                    \right]_{a}
            \end{array}
            \right]_{i}
    }
$$
\[
    \left[
        \begin{array}{c}
            \left[
                \begin{array}{ccc}
                    [1\ 2\ 3\ 4]_{b} & [5\ 6\ 7\ 8]_{b} & [9\ 10\ 11\ 12]_{b}
                \end{array}
                \right]_{a}
            \\
            \left[
                \begin{array}{ccc}
                    [13\ 14\ 15\ 16]_{b} & [17\ 18\ 19\ 20]_{b} & [21\ 22\ 23\ 24]_{b}
                \end{array}
                \right]_{a}
        \end{array}
        \right]_{i}
\]

\[
    \left[
        \begin{array}{c}
            \left[
                \begin{array}{ccc}
                    \left[\begin{array}{rrrr}
                                  \phantom{0}1 & \phantom{0}2 &
                                  \phantom{0}3 & \phantom{0}4
                              \end{array}\right]_{b}
                     &
                    \left[\begin{array}{rrrr}
                                  \phantom{0}5 & \phantom{0}6 &
                                  \phantom{0}7 & \phantom{0}8
                              \end{array}\right]_{b}
                     &
                    \left[\begin{array}{rrrr}
                                  \phantom{0}9 & 10 & 11 & 12
                              \end{array}\right]_{b}
                \end{array}
                \right]_{a}
            \\[12pt]
            \left[
                \begin{array}{ccc}
                    \left[\begin{array}{rrrr}
                                  13 & 14 & 15 & 16
                              \end{array}\right]_{b}
                     &
                    \left[\begin{array}{rrrr}
                                  17 & 18 & 19 & 20
                              \end{array}\right]_{b}
                     &
                    \left[\begin{array}{rrrr}
                                  21 & 22 & 23 & 24
                              \end{array}\right]_{b}
                \end{array}
                \right]_{a}
        \end{array}
        \right]_{i}
\]

TMP




\section{Reference}
\bibliographystyle{plain}  % or another style like unsrt, IEEEtran, etc.
\bibliography{refer}  % references.bib is the file name


% \subsection{Definition of EMH}
% \begin{definition}[EMH means Symmetries]
%     \label{def:mydefinition}
%     We say a market is completely EMH if and
%     only if for every condition $A$, the random
%     variable of the logarithmic growth rate $W$
%     agrees that $\mathbb{E}_{W\mid A}(W)\doteq 0$.
% \end{definition}


% \subsection{Prior Work}
% Discussion of what others have done.




% \section{Methods}
% How you approached the problem.




% \section{Results}
% What you found.




% \section{Conclusion}
% Your closing thoughts.




